\documentclass[11pt]{article}
\usepackage{palestra}
\clavetrue
% la clave se consulta, no se lee de corrido: compacta
\newcommand{\cara}[1]{\par\vspace{1.2mm}{\rotulo\footnotesize\color{oro}#1}\par\vspace{.6mm}}
\newcommand{\nj}[1]{\textbf{#1.}\ }
\newcommand{\nota}[1]{{\itshape\color{tinta!75} #1}}
\newcommand{\sep}{\quad·\quad}

\begin{document}\small\setlength{\parskip}{1mm}

% ═════════════════════════════ FORMATO DE 2.º ═════════════════════════════
\begin{ficha}{Clave}{Clave · Ficha 1 · Fuerza y bando}{1}
\cara{DELANTE · EL NÚMERO Y SU SIGNO}
\nj{1} a) $-3$, 3, ←\sep b) $+5$, 5, →\sep c) $-1$, 1, ←\sep d) $+2$, 2, →\sep
e) $-7$, 7, ←\sep f) $+4$, 4, →. \nota{La fuerza va en la cabeza del tirador; el
signo lo da el color.}

\nj{2} a) $-6$: $|{-6}| = 6 > 4$\sep b) $+5$: $5 > 2$\sep c) $-7$: $7 > 3$.

\nj{3} $-8 < -5 < -2 < 0 < 3 < 7$.\sep a) $-3 < +1$\sep b) $-7 < -2$\sep
c) $0 > -5$\sep d) $-1 > -9$. \nota{b) y d) son los de dos negativos: ahí se ve
quién confunde fuerza con orden.}

\textbf{Reto.} El 0; no tira hacia ningún lado, no es de ningún bando. Fuerza 5
tienen dos, $-5$ y $+5$ (opuestos); fuerza 0, solo el 0.

\cara{REVERSO · MÁS DUELOS}
\nj{1} a) $-8$; $+3$; no\sep b) $-5$; $-2$; no\sep c) $-4$; $+1$; no\sep
d) $-7$; $-1$; no\sep e) $+6$; $+6$; sí. Gana el mismo las dos cuando el que tira
más fuerte es \textbf{positivo}.

\nj{2} Moscú $-8$ < Oslo $-6$ < Burgos $-4$ < Madrid $+1$ < Sevilla $+6$. Más
lejos del 0, Moscú: la más fría.

\nj{3} a) El buzo ($12 > 8$); la gaviota\sep b) el garaje; 2 plantas\sep
c) $-325$ y $-287$; Euclides ($-325 < -287$).

\textbf{Reto.} En el bando azul, sí: $+7 > +2$. En el rojo, no: $-7$ tira más que
$-2$ y es menor.
\end{ficha}

\vspace{3mm}
\begin{ficha}{Clave}{Clave · Ficha 2 · Sumar y restar}{2}
\cara{DELANTE · ENTRAR Y RETIRARSE}
\nj{1} a) $-3 + 5 = +2$\sep b) $-4 + (-2) = -6$\sep c) $-6 + 2 + 3 = -1$.
\nota{Vale cualquier orden de los que entran: la bandera acaba igual.}

\nj{2} a) sí; nadie; $+3$\sep b) no; $-4$ y $+4$; $+5$\sep c) no; $+3$ y $-3$; $-5$\sep
d) sí; nadie; 0.

\nj{3} a) 3\sep b) $-7$\sep c) 8\sep d) 0\sep e) $-2$\sep f) $-8$\sep g) $-5$\sep h) 3.

\textbf{Reto.} Sí: restando un negativo, $2 - (-3) = 5 > 2$. Y sumando un negativo
se acaba más abajo: $2 + (-3) = -1 < 2$.

\cara{REVERSO · MÁS JUGADAS}
\nj{1} a) $-4 + 6 - (-4) = +6$\sep b) $3 - 5 = -2$\sep c) $-2 + (-2) - (-5) = +1$.
En c) entra el $-5$ con su pareja, el $+5$. \nota{En b) también vale $3 - (+5)$.}

\nj{2} a) de $-2$, 5 a la derecha: $+3$\sep b) de 3, 7 a la izquierda: $-4$\sep
c) de $-4$, 6 a la derecha: $+2$. Hacia la derecha: se va uno que tiraba hacia la
izquierda.

\nj{3} a) $-5 + 9 = +4$\,°C\sep b) $30 - 45 = -15$\,€\sep c) $-212 - (-287) = 75$
años\sep d) $-40 + 15 - 30 = -55$\,m.

\textbf{Reto.} Por ejemplo $-3 + 1 + 2 = 0$; con resta, $5 - 2 - 3 = 0$.
\end{ficha}

\newpage
\begin{ficha}{Clave}{Clave · Ficha 3 · Las pociones}{3}
\cara{DELANTE · MULTIPLICAR Y DIVIDIR}
\nj{1} a) $+6$\sep b) $-8$\sep c) $+3$\sep d) $+3$\sep e) $-6$\sep f) $+4$. Cambian
de bando los de calavera; los verdes y los azules lo dejan en el suyo.

\nj{2} a) $(-2)\cdot 3 = -6$; $-6$; $-1$\sep b) $(-3)\cdot(-2) = +6$; $+6$; $+2$\sep
c) $(+8)\dv 2 = +4$; $+4$; $-3$\sep d) $2\cdot(-4) = -8$; $-8$; $-5$.

\nj{3} a) 12\sep b) $-4$\sep c) $-20$\sep d) 5\sep e) $-1$\sep f) 12\sep g) $-3$\sep h) $-12$.

\textbf{Reto.} De $+6$ a $-2$: dividir entre $-3$ (en la palestra no hay ese
frasco: un $\dv 3$ y una traición). De $-6$ a $+12$: $\cdot(-2)$. Para dejarlo
igual, $\cdot 1$, que no está; o dos traiciones seguidas.

\cara{REVERSO · EL PARÉNTESIS}
\nj{1} a) $+3$; $\cdot 3$; $+9$; $+9$\sep b) $-3$; $\cdot(-2)$; $+6$; $+6$\sep
c) $-5$; el menos; $+5$; $+11$\sep d) $-3$; $\cdot 2$; $-6$; $-1$. El menos de
delante le cambia de bando, como un frasco de $\cdot(-1)$.

\nj{2} a) $4 + 6 = 10$\sep b) $-2\cdot(-2) + 1 = 5$\sep c) $3 - 3 = 0$\sep
d) $5 - 6 = -1$\sep e) $4\dv(-2) = -2$\sep f) $5 - 2\cdot(-3) = 11$.

\nj{3} a) $(-3)\cdot 4 = -12$\,m\sep b) $5 + 4\cdot(-2) = -3$\,°C\sep
c) $(-18)\dv 3 = -6$\,€.

\textbf{Reto.} No: $-(2 - 5) = -(-3) = +3$ y $-2 - 5 = -7$. \nota{En el campo: en la
primera, el paréntesis se reduce en el banquillo y el menos le cambia de bando.}
\end{ficha}

\vspace{3mm}
\begin{ficha}{Clave}{Clave · Ficha 4 · Potencias}{4}
\cara{DELANTE · SU PROPIO FRASCO}
\nj{1} a) $+4$\sep b) $+9$\sep c) $+9$\sep d) $+8$\sep e) $-8$. En e), tras el
primer frasco, $+4$ (azul); tras el segundo, $-8$ (rojo).

\nj{2} a) $9 > -9$\sep b) $-8 = -8$\sep c) $1 > -1$\sep d) $4 = 4$. Lo mismo en b),
porque el exponente es impar, y en d), porque un número y su opuesto tienen el
mismo cuadrado.

\nj{3} a) $5 - 4 = 1$\sep b) $9 - 8 = 1$\sep c) $-4 + 5 = 1$\sep d) $4 - 8 = -4$\sep
e) $-1 - 9 = -10$\sep f) $10 - 12 = -2$.

\textbf{Reto.} $(-1)^{\text{par}} = 1$ y $(-1)^{\text{impar}} = -1$:
$(-1)^{100} = 1$, $(-1)^{101} = -1$. Cada dos traiciones, vuelve a su bando.

\cara{REVERSO · EL OJO DE ARGOS}
\nj{1} a) la 2.ª: es $-8 + 6$, y da $-2$\sep b) la 2.ª: $(-2)^3 = -8$, y da $-3$\sep
c) la 2.ª: el paréntesis vale 3, y da $-3$\sep d) la 3.ª: $4 - (+4) = 0$.
\nota{a) y d) son el mismo fallo: confundir retirar un negativo con retirar un
positivo.}

\nj{2} a) $18 + 5 = 23$\sep b) $9 - 10 = -1$\sep c) $-4 + 8 = 4$\sep
d) $4\cdot(-3) + 1 = -11$\sep e) $-18\dv 9 = -2$\sep f) $-8 - 1 = -9$.

\nj{3} Abierto. Por ejemplo, $(-2)^2\cdot(-1) + 1 = -3$.

\textbf{Reto.} $-2^2$ sale a tirar con $-4$; $(-2)^2$, con $+4$.
\end{ficha}

% ═════════════════════════════ FORMATO DE 1.º ═════════════════════════════
\newpage
\begin{ficha}{Clave}{Clave · Formato de 1.º}{0}
\cara{FICHA 1 · FUERZA Y BANDO}
\textbf{Delante.} 2) a) $+3$, 3\sep b) $-2$, 2\sep c) $-6$, 6\sep d) $+5$, 5.
3) a) $-3 < +2$\sep b) $-5 < -1$\sep c) $+4 > -6$. Reto: el 0, que no es de ningún
bando.
\textbf{Reverso.} 1) a) $-6$\sep b) $-7$\sep c) $-8$\sep d) $+5$. 2) a) $<$\sep b) $>$\sep
c) $>$\sep d) $>$. 3) a) en Burgos\sep b) el trastero.

\cara{FICHA 2 · SUMAR Y RESTAR}
\textbf{Delante.} 2) a) $-4 + 6 = +2$\sep b) $-2 + (-3) = -5$. 3) a) $-3$\sep b) $-4$\sep
c) $-5$\sep d) 0. Reto: dos opuestos; tantas maneras como fuerzas ($-1$ y $+1$,
$-2$ y $+2$…).
\textbf{Reverso.} 2) a) sí; nadie; $+3$\sep b) no; $-2$ y $+2$; $+3$\sep c) no; $+2$ y
$-2$; $-5$. 3) a) 6\sep b) $-6$.

\cara{FICHA 3 · LAS POCIONES}
\textbf{Delante.} 2) a) $+8$\sep b) $-6$\sep c) $+5$\sep d) $-6$. 3) a) 8\sep b) $-2$\sep
c) $-9$\sep d) 2. Reto: $\cdot(-1)$, la traición.
\textbf{Reverso.} 2) a) $3\cdot(-2) = -6$; $-2$\sep b) $(-2)\cdot(-3) = +6$; $+5$\sep
c) $(+6)\dv 2 = +3$; $-2$. 3) a) $(-2)\cdot 5 = -10$\,m\sep b) $(-12)\dv 3 = -4$\,€.

\cara{FICHA 4 · POTENCIAS}
\textbf{Delante.} 2) a) $+4$\sep b) $+9$\sep c) $+1$\sep d) $-8$. 3) a) 4 y $-4$\sep
b) 1 y $-1$. Reto: $-1$, $+1$, $+1$.
\textbf{Reverso.} 2) a) 7\sep b) 5\sep c) $-7$\sep d) 2. 3) a) la 1.ª: $(-3)^2 = 9$, y da
10\sep b) la 2.ª: $2 - 1 = 1$.
\end{ficha}
\end{document}
